\begin{center}
    {\LARGE \textbf{2014分析化学}} \\[2ex]
    {\large 时量180分钟 \quad 满分90分}
\end{center}

\vspace{0.5cm}
\noindent\textbf{注意事项:}
\begin{enumerate}[label=\arabic*., parsep=0pt, itemsep=0pt]
    \item 答题前,考生先将自己的姓名、学号、考场号、座位号填写在答题卡上,将条形码准确粘贴在考生信息条形码粘贴区。
    \item 考试结束后,将本试卷与答题纸一并收回。
    \item 回答各个问题时,将答案写在答题纸上,写在本试卷上无效。
\end{enumerate}

\vspace{0.5cm}
\noindent\textbf{一、简答题(35 pts.)}

\begin{enumerate}[label=\arabic*., itemsep=1.5ex]
    \item 碘量法定铜的主要误差来源是什么？加入$\ce{NH4SCN}$的作用是什么？
    \item 在弱碱性溶液中用EDTA滴定$\ce{Zn^{2+}}$常使用$\ce{NH3-NH4+}$溶液，其作用是什么？
    \item 如配制EDTA的水中含有少量的$\ce{Zn^{2+}}$和$\ce{Ca^{2+}}$，在$\mathrm{pH}=5.0$时用标准$\ce{Zn^{2+}}$标定EDTA，在$\mathrm{pH}=2.0$时测定$\ce{Fe^{3+}}$，结果如何？在$\mathrm{pH}=5.0$时测$\ce{Pb^{2+}}$结果如何？在$\mathrm{pH}=10$时测$\ce{Mg^{2+}}$和$\ce{Ca^{2+}}$合量，结果如何？
    \item 饮用水中含有少量$\ce{CHCl3}$，取水样100mL，用10mL戊醇萃取，有90.5\%的$\ce{CHCl3}$被萃取，计算取10mL水样用10mL戊醇萃取时，$\ce{CHCl3}$被萃取的百分率。
    \item 用一般分光光度法测量$0.00100\mathrm{mol/L}$锌标准溶液和含锌试液，分别测得$A=0.700$和
    \\
    $A=1.000$，两种溶液的透光率相差多少？如果用$0.00100\mathrm{mol/L}$锌标准溶液做参比溶液，试液的吸光度又是多少？示差分光光度法与一般分光光度法相比较，读数标尺放大了多少倍？
    \item 以$0.10\mathrm{mol/L}\ \ce{Fe^{3+}}$的$\ce{HCl}$液滴定$20.00\mathrm{mL}\ 0.05\mathrm{mol/L}$的$\ce{SnCl2}$溶液，计算化学计量点及计量点前后$0.1\%$的电位值。(已知$E^{\ominus\prime}_{\ce{Fe^{3+}/Fe^{2+}}}=0.68\mathrm{V}$，$E^{\ominus\prime}_{\ce{Sn^{4+}/Sn^{2+}}}=0.14\mathrm{V}$)
    \item 有一单色光通过厚度为1cm的有色溶液，其强度减弱20\%。若通过5cm厚度的相同溶液，其强度减弱百分之几？
\end{enumerate}

\vspace{0.5cm}
\noindent\textbf{二、计算题(55 pts. total)}

\begin{enumerate}[label=\arabic*., itemsep=1.5ex]
    \item (12 pts.) 计算$\ce{Ag2S}$在氨和氯化铵浓度分别为$0.20\mathrm{mol/L}$和$0.10\mathrm{mol/L}$溶液中的溶解度。(已知$\ce{Ag2S}$的$K_{\mathrm{sp}}=2\times10^{-49}$，$\ce{NH3}$的$K_b=1.8\times10^{-5}$，$\ce{Ag^+ - NH3}$络合物的$\lg\beta_1,\lg\beta_2$分别为$3.24$，$7.05$；$\ce{H2S}$的$K_{a1}=1.3\times10^{-7}$，$K_{a2}=7.1\times10^{-15}$)
    \item (12 pts.) 在$\mathrm{pH}=8.0$的氨性缓冲溶液中，用$1.0\times10^{-2}\mathrm{mol/L}$EDTA溶液滴定等浓度的$\ce{Zn^{2+}}$，终点时游离氨的浓度为$0.20\mathrm{mol/L}$；在此条件下能否准确滴定$\ce{Zn^{2+}}$ ($E_t \leq 0.1\%$)？如果要求$E_t<0.3\%$，则要求$\Delta\mathrm{pZn}'$为多少？（已知$\mathrm{pH}=8.0$时$\lg\alpha_{Y(\mathrm{H})}=0.45$；$\ce{Zn^{2+}-NH3}$络合物的$\lg\beta_1-\lg\beta_4$分别为$2.37$，$4.81$，$7.31$，$9.46$；$\lg K_{\ce{ZnY}}=16.50$）
    \item (10 pts.) 忽略离子强度的影响，计算$\mathrm{pH}=9.00$，$c_{\ce{NH3}}=0.10\mathrm{mol/L}$溶液中，$\ce{Zn^{2+}/Zn}$电对的条件电极电位。(已知$\ce{NH3}$的$\mathrm{p}K_b=4.74$；$\ce{Zn^{2+}-NH3}$络合物的$\lg\beta_1-\lg\beta_4$分别为$2.37$，$4.81$，$7.31$，$9.46$；$E^\ominus_{\ce{Zn^{2+}/Zn}}=-0.763\mathrm{V}$)
    \item (11 pts.) 某制药厂生产维生素B1片剂，要求每10g药片中杂质铁控制在128mg。从某批产品中取样(10g)进行5次分析，结果(mg)为127，128，128，130，131，求相对标准差及置信度95\%的置信区间，这批产品杂质Fe含量是否超过规定标准(95\%置信水平)？
    \\
    已知：
    \begin{center}
    \begin{tabular}{l|ccc}
    $f$ & 3 & 4 & 5 \\
    \hline
    $t_{0.05}$ & 3.18 & 2.78 & 2.57 \\
    $t_{0.10}$ & 2.35 & 2.13 & 2.02 \\
    \end{tabular}
    \end{center}
    \item (10 pts.) 用$0.1000\mathrm{mol/L}\ \ce{NaOH}$溶液滴定$0.100\mathrm{mol/L}$羟胺盐酸盐($\ce{NH3^+OH\cdot Cl^-}$)
    \\
    和$0.100\mathrm{mol/L}\ \ce{NH4Cl}$混合溶液，问：
    \begin{enumerate}[label=(\arabic*), noitemsep, topsep=0pt, partopsep=0pt]
        \item 化学计量点时溶液的pH值为多少？
        \item 化学计量点时有百分之几的$\ce{NH4Cl}$参加了反应？(已知$\ce{NH2OH}$的$\mathrm{p}K_b=9.00$，$\ce{NH3}$的$\mathrm{p}K_b=4.74$)
    \end{enumerate}
\end{enumerate}